\documentclass[11pt,a4paper]{article}

\usepackage[margin=2.5cm]{geometry}
\usepackage[T1]{fontenc}
\usepackage[utf8]{inputenc}
\usepackage{lmodern}
\usepackage{microtype}
\usepackage{amsmath,amssymb,mathtools}
\usepackage{booktabs}
\usepackage{graphicx}
\usepackage{siunitx}
\usepackage{enumitem}
\usepackage{xcolor}
\usepackage{hyperref}
\usepackage[capitalise,noabbrev]{cleveref}
\usepackage[numbers,sort&compress]{natbib}

\hypersetup{
  colorlinks=true,
  linkcolor=blue!50!black,
  citecolor=blue!50!black,
  urlcolor=blue!50!black
}

\newcommand{\R}{\mathbb{R}}
\newcommand{\E}{\mathbb{E}}
\newcommand{\abs}[1]{\left|#1\right|}
\newcommand{\norm}[1]{\left\lVert#1\right\rVert}
\newcommand{\sign}{\operatorname{sign}}
\newcommand{\todo}[1]{\textcolor{red}{[TODO: #1]}}

\title{Neuromorphic Federated Learning via Leaky Integrate-and-Fire Gradient Communication\\
\large Living Technical Report and Project Knowledge Base}
\author{NeuromorphicFL Project}
\date{\today}

\begin{document}
\maketitle

\begin{abstract}
This report is a living technical knowledge base for a research direction that transfers the integrate--threshold--spike principle from neuromorphic control to distributed and federated optimization. The central idea is not to train a spiking neural network. Instead, ordinary stochastic gradients are treated as continuous or discrete evidence signals that drive leaky integrate-and-fire (LIF) communication states. Communication occurs only when sufficient signed evidence has accumulated, and the global model evolves through sparse fixed-amplitude signed parameter-update events. The event timing or rate therefore carries gradient-magnitude information, while the sign carries update direction. The project is motivated by the hybrid-systems viewpoint developed for neuromorphic control, especially the work of Petri, Scheres, Steur, and Heemels, but must be distinguished carefully from established work on one-bit gradients, error feedback, gradient sparsification, lazy/event-triggered distributed learning, and event-based optimization. This document consolidates the current mathematical formulation, literature position, first scalar stochastic experiments, observed strengths and limitations, and the roadmap toward an asynchronous federated algorithm.
\end{abstract}

\tableofcontents
\newpage

\section{Motivation and Research Question}

\subsection{From neuromorphic control to neuromorphic communication}

Recent work on neuromorphic control develops feedback controllers whose information and actuation are represented by sparse fixed-amplitude spikes rather than continuously varying control values. In the simple two-neuron setting of \citet{petri2024simple}, leaky integrate-and-fire (LIF) states accumulate plant-output information, fire when a threshold is reached, and generate impulsive control actions. The resulting closed loop is naturally represented as a hybrid dynamical system. The work emphasizes practical stability, dwell-time guarantees, and explicit relationships between neuron parameters and closed-loop behavior. The later emulation-based formulation \citep{petri2025emulation} asks whether a spike train can emulate a continuous stabilizing input closely enough that stability properties can be transferred.

The present project asks whether the same structural principle can be applied not to physical actuation but to \emph{learning communication}. The desired analogy is
\begin{center}
\begin{adjustbox}{max width=\textwidth}
\begin{tabular}{ll}
\toprule
Neuromorphic control & Neuromorphic distributed learning \\
\midrule
plant/control error & stochastic gradient information \\
LIF membrane state & communication evidence state \\
threshold crossing & communication event \\
fixed signed control spike & fixed signed parameter-update spike \\
spike timing/rate & gradient-magnitude encoding \\
closed-loop state & model parameter vector \\
Lyapunov stability & optimization convergence/descent \\
\bottomrule
\end{tabular}
\end{adjustbox}
\end{center}

The intended neural network does \emph{not} need to be a spiking neural network. Backpropagation can remain completely conventional. The candidate novelty is an event-driven communication and optimization layer that transforms local stochastic-gradient information into sparse signed events.

\subsection{Core research question}

The working research question is
\begin{quote}
Can stochastic learning information be accumulated locally by leaky integrate-and-fire dynamics and communicated as sparse fixed-amplitude events while preserving useful optimization convergence properties and, in a federated setting, reducing communication without sacrificing learning quality?
\end{quote}

This question contains several subproblems that should be kept separate:
\begin{enumerate}[label=(\roman*)]
  \item \textbf{Temporal encoding:} Can gradient magnitude be represented through spike timing or rate rather than transmitted amplitude?
  \item \textbf{Statistical reliability:} Does temporal evidence accumulation reduce wrong-sign updates under stochastic gradients?
  \item \textbf{Communication efficiency:} Does the event mechanism reduce the number of messages and ultimately transmitted bits at a given target error?
  \item \textbf{Memory/freshness:} What role should leakage play when accumulated evidence becomes stale because the optimization landscape or global model changes?
  \item \textbf{Federation:} Can multiple clients asynchronously generate events from heterogeneous local objectives while retaining convergence of a global model?
\end{enumerate}

\subsection{What the project is not}

Several nearby formulations are intentionally excluded from the core claim.
\begin{itemize}
  \item The project is not simply ``federated learning of an SNN.'' Federated SNN training is a separate established topic.
  \item The project is not simply one-bit gradient communication. One-bit SGD and sign-based distributed optimization are established \citep{seide2014onebit,bernstein2018signsgd}.
  \item The project is not simply thresholding accumulated compression residuals. Error feedback, sparsified SGD with memory, and deep gradient compression already preserve discarded gradient information across iterations \citep{lin2017dgc,stich2018memory,richtarik2021ef21}.
  \item The project is not simply event-triggered communication. Lazy and event-based distributed optimization already reduce communication by skipping uninformative exchanges \citep{chen2018lag,er2025eventbased}.
\end{itemize}

The research direction is therefore only compelling if the neuromorphic viewpoint produces a \emph{coherent hybrid optimization architecture}: a stateful temporal encoder that determines when to communicate, represents update magnitude through event timing/rate, supports explicit memory/freshness dynamics, and admits systems-style convergence and inter-event analysis.

\subsection{Long-term federated interpretation}

For a federated objective
\begin{equation}
  F(w)=\sum_{i=1}^N p_i F_i(w),
\end{equation}
each client will eventually maintain a communication state $z_i$ driven by its local stochastic gradients. The global model should change only when events arrive. In the strongest version of the architecture there is no fixed concept of a federated ``round.'' Instead, clients compute locally at their own pace, accumulate evidence, and emit sparse events only when their local communication states cross thresholds. This would turn the full learning process into an interconnected hybrid dynamical system.

\section{Literature Position and Novelty Constraints}

\subsection{Federated learning and communication efficiency}

Federated learning (FL) trains a shared model from decentralized data by aggregating locally computed updates rather than centralizing the raw data. Federated averaging (FedAvg) established the canonical round-based architecture and already identified communication as a principal systems constraint \citep{mcmahan2017fedavg}. The present project is therefore situated within a mature communication-efficient learning literature and must be evaluated against methods that reduce message precision, message sparsity, or message frequency.

\subsection{One-bit and sign-based optimization}

One-bit SGD showed that aggressive gradient quantization can be practical if quantization error is carried forward across minibatches \citep{seide2014onebit}. signSGD later formalized sign-only stochastic optimization and its distributed majority-vote variant \citep{bernstein2018signsgd}. These works already establish that fixed signed messages can be useful. Consequently, transmitting only $\pm1$ is not itself novel.

The proposed LIF mechanism differs in how magnitude information is represented. Periodic signSGD acts on the instantaneous sign
\begin{equation}
  w_{k+1}=w_k-q\,\sign(g_k),
\end{equation}
whereas a neuromorphic temporal encoder accumulates gradient evidence and communicates only after a threshold crossing. The magnitude of a persistent gradient can therefore be represented through the inter-event interval or firing rate.

\subsection{Error feedback, sparsification, and residual memory}

Deep Gradient Compression (DGC) substantially reduces distributed gradient traffic by sparsifying communication and retaining information needed to preserve learning quality \citep{lin2017dgc}. Sparsified SGD with Memory gives a theoretical treatment of compression with accumulated error compensation \citep{stich2018memory}. EF21 develops a modern error-feedback mechanism with convergence guarantees for heterogeneous distributed optimization \citep{richtarik2021ef21}. Sparse Ternary Compression (STC) targets communication-efficient FL under non-IID data \citep{sattler2019stc}.

These methods are a critical novelty constraint. A zero-leak accumulator of the form
\begin{equation}
  z_{k+1}=z_k-\gamma g_k
\end{equation}
followed by a quantization/threshold operation is closely related to error-feedback and sigma--delta ideas. Therefore, the project must not claim that ``accumulate small gradients until they are large enough'' is new.

The candidate differentiators are instead:
\begin{enumerate}[label=(\alph*)]
  \item a temporal rate code in which gradient magnitude determines event timing;
  \item an explicit LIF state with a tunable forgetting timescale;
  \item a hybrid-system flow/jump model of optimization;
  \item potentially Lyapunov/descent-derived event conditions rather than heuristic thresholds;
  \item asynchronous bidirectional event representations of global model evolution.
\end{enumerate}

\subsection{Lazy and event-based distributed learning}

LAG reduces communication by reusing stale gradients when local gradients change slowly and establishes convergence properties under standard optimization assumptions \citep{chen2018lag}. More recent event-based distributed optimization directly makes communication conditional on local events. In particular, \citet{er2025eventbased} develop an event-based ADMM framework with convergence analysis under heterogeneous local data distributions and demonstrate distributed learning experiments.

Thus, event-triggered communication is also not new. The useful research question is whether the \emph{neuromorphic encoding mechanism itself} provides a new optimization object: local evidence states that integrate stochastic gradients, fire signed fixed-amplitude model updates, and jointly determine both communication frequency and update information.

\subsection{Relationship to neuromorphic control}

The control-theoretic inspiration is strongest in the hybrid model of \citet{petri2024simple}. There, neuron-like states integrate measured plant information, threshold crossings create spikes, and the controller is analyzed through continuous flows and discrete jumps. Practical stability and strictly positive dwell times are central quantities. The emulation-based work \citep{petri2025emulation} is especially relevant conceptually because it starts from a desired continuous signal and designs a spiking representation whose approximation properties are then connected to closed-loop stability.

The optimization analogue is to start from ideal gradient flow
\begin{equation}
  \dot w=-\kappa \nabla F(w)
\end{equation}
and construct a spiking parameter trajectory that emulates the accumulated gradient action with bounded communication-state error.

\subsection{Current novelty statement: deliberately provisional}

The literature reviewed so far supports only the following cautious statement:
\begin{quote}
A candidate research gap exists in formulating federated/distributed stochastic optimization as a hybrid dynamical system whose local stochastic gradients drive LIF-like communication states and whose global model evolves through asynchronous fixed signed parameter-update events. The potentially distinctive contribution lies in the temporal evidence code, leak-controlled information freshness, hybrid convergence analysis, and event-driven model trajectory---not in gradient accumulation, compression, sign messages, or event triggering individually.
\end{quote}

This must be stress-tested continuously as the literature review deepens. Before any paper-level novelty claim, the project should explicitly compare against error-feedback compressors, sign-based optimization, LAG/event-triggered methods, and event-based distributed optimization.

\section{Candidate LIF-SGD Formulation}

\subsection{Scalar quadratic starting point}

The diagnostic problem is
\begin{equation}
  F(w)=\frac{a}{2}(w-\theta)^2,\qquad a>0,
\end{equation}
with gradient $\nabla F(w)=a(w-\theta)$. Ideal continuous gradient flow is
\begin{equation}
  \dot w=-\kappa a(w-\theta),
\end{equation}
which converges exponentially to $w^\star=\theta$.

\subsection{Zero-leak integrate-and-fire encoder}

Introduce a communication state $z$. Between events the model is held fixed and the desired gradient action is accumulated:
\begin{align}
  \dot w &= 0,\\
  \dot z &= -\kappa a(w-\theta).
\end{align}
A communication event occurs when $|z|=\Delta$. The model then jumps by a fixed quantum $q$,
\begin{align}
  s &= \sign(z),\\
  w^+ &= w+q s,\\
  z^+ &= 0,
\end{align}
for the ideal continuous threshold-crossing model. In discrete implementations a subtractive reset $z^+=z-s\Delta$ is preferable because it preserves threshold overshoot.

For fixed $w\neq\theta$ and zero leak, the inter-event interval is
\begin{equation}
  \tau=\frac{\Delta}{\kappa |\nabla F(w)|},
\end{equation}
so that
\begin{equation}
  f=\frac{1}{\tau}=\frac{\kappa |\nabla F(w)|}{\Delta}.
\end{equation}
Thus gradient magnitude is represented by event rate rather than event amplitude.

If $q=\Delta$ and events are indexed by $j$, the deterministic scalar event map becomes
\begin{equation}
  w_{j+1}=w_j-q\,\sign(w_j-\theta),
\end{equation}
which is identical to constant-step sign descent in event index. This is an important negative result: the deterministic fixed-threshold jump sequence is not by itself a novel optimizer. Its distinctive information is the endogenous timing of those jumps.

For fixed $q$, the event-indexed trajectory enters a practical neighborhood of the optimum. In the scalar deterministic case,
\begin{equation}
  \limsup_{t\rightarrow\infty}|w(t)-\theta|<q
\end{equation}
under the simple fixed-amplitude construction. Exact convergence requires a diminishing jump, adaptive quantization, or another refinement.

\subsection{Stochastic gradients and LIF dynamics}

The main candidate mechanism uses stochastic gradients
\begin{equation}
  g_k=a(w_k-\theta)+\epsilon_k,
\end{equation}
with zero-mean noise $\epsilon_k$. The discrete LIF state is
\begin{equation}
  z_{k+1}=\rho z_k-\gamma g_k,
  \qquad 0\leq\rho\leq1.
\end{equation}
The relation to a continuous leak rate $\mu$ over sampling interval $h$ is approximately $\rho=e^{-\mu h}$. A threshold crossing $|z|\geq\Delta$ creates event $s=\sign(z)$ and model update
\begin{equation}
  w^+=w+q s,
\end{equation}
followed by subtractive reset
\begin{equation}
  z^+=z-s\Delta.
\end{equation}
The three parameters have deliberately separate roles:
\begin{center}
\begin{tabular}{cl}
\toprule
$q$ & optimization/update resolution \\
$\Delta$ & amount of temporal evidence required before communication \\
$\rho$ (or $\mu$) & memory/freshness timescale \\
\bottomrule
\end{tabular}
\end{center}

Before firing, the state is an exponentially weighted accumulation of past gradients,
\begin{equation}
  z_k\approx-\gamma\sum_{r=0}^{\infty}\rho^r g_{k-r}.
\end{equation}
Hence the event sign is not the sign of a single noisy minibatch gradient. It is the sign of accumulated evidence after enough magnitude has built up to cross a threshold.

\subsection{Continuous stochastic interpretation}

A diffusion approximation is
\begin{equation}
  dz=(-\mu z-\kappa a(w-\theta))dt-\kappa\sigma dB_t.
\end{equation}
For fixed $w$ and ignoring threshold/reset, this is an Ornstein--Uhlenbeck process with mean
\begin{equation}
  \E[z]=-\frac{\kappa a(w-\theta)}{\mu}
\end{equation}
and stationary variance
\begin{equation}
  \operatorname{Var}(z)=\frac{\kappa^2\sigma^2}{2\mu}.
\end{equation}
Without noise, a leaky state fires only if
\begin{equation}
  |\nabla F(w)|>\frac{\mu\Delta}{\kappa},
\end{equation}
which creates a deterministic gradient deadzone. This suppresses weak/noisy communication but can also slow learning.

For zero leak, the Brownian first-passage model yields a wrong-boundary probability of the form
\begin{equation}
  p_{\mathrm{wrong}}(r)=\frac{1}{1+\exp\!\left(\frac{2ar\Delta}{\kappa\sigma^2}\right)},
  \qquad r=|w-\theta|,
\end{equation}
showing explicitly that demanding more accumulated evidence improves sign reliability. The corresponding expected signed spike imbalance recovers the underlying gradient direction in expectation. This motivates viewing zero-leak IF as an information-conserving temporal encoder, while leakage intentionally trades information conservation for freshness and reduced noise chatter.

\subsection{Proposed federated lift}

For client $i$ with stochastic local gradient $g_i$, a first asynchronous federated form is
\begin{equation}
  z_i\leftarrow\rho_i z_i-\gamma_i g_i.
\end{equation}
Whenever coordinate $j$ crosses its threshold, the client emits only an event $(j,s_{i,j})$. The server applies
\begin{equation}
  w_j^+=w_j+q_j s_{i,j}
\end{equation}
and may broadcast the accepted event so all clients reconstruct the same global model trajectory. This removes the requirement for fixed FL rounds. The resulting system is an interconnected stochastic hybrid dynamical system and is the long-term theoretical target of the project.

\section{Current Diagnostic Experiments}

\subsection{Experiment 1: stationary stochastic quadratic}

\paragraph{Purpose.}
The first experiment tests whether the LIF communication primitive behaves as intended before introducing federation or neural networks. The objective is
\begin{equation}
  F(w)=\frac12 w^2,
\end{equation}
with stochastic gradients
\begin{equation}
  g_k=w_k+\epsilon_k,\qquad \epsilon_k\sim\mathcal{N}(0,2^2).
\end{equation}
The initial parameter is $w_0=3$. The main run uses 5000 gradient evaluations and 200 Monte Carlo seeds. The LIF parameters are
\begin{equation}
  \gamma=0.05,\qquad \Delta=0.5,\qquad q=0.05,\qquad \rho=0.98.
\end{equation}
For SGD, the learning rate is chosen as
\begin{equation}
  \eta_{\mathrm{SGD}}=\frac{q\gamma}{\Delta}=0.005,
\end{equation}
so that its mean small-signal gradient-flow scale is comparable to the event encoder. signSGD uses the same fixed update quantum $q$ as one IF/LIF event but updates every gradient evaluation.

\paragraph{Baselines.}
The methods are ordinary SGD, periodic signSGD, non-leaky IF-SGD ($\rho=1$), and LIF-SGD. The experiment additionally holds $w$ fixed at several values to measure communication-event rate and the probability that an emitted signed event points in the wrong direction. A leak sweep varies $\rho\in\{1,0.995,0.99,0.98,0.95,0.90\}$.

\paragraph{Observed optimization behavior.}
For the current configuration, the 200-seed averages were approximately
\begin{center}
\begin{adjustbox}{max width=\textwidth}
\begin{tabular}{lrrrr}
\toprule
Method & final $\E|w|$ & tail $\E[w^2]$ & mean $w^2$ & communications \\
\midrule
SGD & 0.0793 & 0.01017 & 0.18973 & 5000 \\
signSGD & 0.2020 & 0.06218 & 0.11658 & 5000 \\
IF-SGD & 0.0790 & 0.01035 & 0.19226 & 206.3 \\
LIF-SGD ($\rho=0.98$) & 0.0743 & 0.00934 & 0.20122 & 177.8 \\
\bottomrule
\end{tabular}
\end{adjustbox}
\end{center}
Thus, in this deliberately small diagnostic, IF/LIF reach an SGD-like tail error with roughly 4\% of the periodic communication count. This number must not be interpreted as an FL bandwidth claim: the problem is scalar, event addressing is absent, and periodic SGD/signSGD are intentionally simple references.

\paragraph{Temporal evidence improves sign reliability.}
The strongest observation is obtained by fixing $w$ while feeding noisy gradients. For example, at $|w|=0.5$, the instantaneous signSGD decision is wrong about 40\% of the time under the selected noise level, whereas the measured LIF wrong-event probability is about 3.5\%. At $|w|=1$, the corresponding values are roughly 30.9\% and 0.14\%. This supports the interpretation that LIF is not merely a low-bit representation: it implements a temporal evidence test before a signed update is emitted.

\paragraph{Rate code.}
The measured event rate increases approximately linearly with $|w|$ in the high-signal regime, consistent with
\begin{equation}
  f\approx\frac{\gamma|\nabla F|}{\Delta}.
\end{equation}
Near the optimum, stochastic noise still creates occasional events. Leakage reduces this chatter.

\paragraph{Leak sweep.}
Stronger leak reduces communication and tail variance but slows the transient because useful evidence is forgotten before it reaches threshold. For the current sweep, communication decreased from roughly 206 events at $\rho=1$ to 97 at $\rho=0.9$, while whole-run mean squared error increased. This establishes the first empirical communication--responsiveness trade-off controlled by the memory parameter.

\subsection{Experiment 2: suddenly moving optimum}

\paragraph{Purpose.}
The second experiment asks whether leakage is beneficial when stored gradient evidence becomes stale. The objective becomes
\begin{equation}
  F_k(w)=\frac12(w-\theta_k)^2,
\end{equation}
with
\begin{equation}
  \theta_k=\begin{cases}
  0,&k<500,\\
  2,&k\geq500.
  \end{cases}
\end{equation}
The experiment uses the same gradient-noise level and LIF parameterization and averages 300 Monte Carlo runs.

\paragraph{Critical plotting convention.}
The primary performance curves must be mean absolute error or RMSE to the instantaneous optimum. The signed ensemble mean $\E[w]$ is not a valid primary convergence diagnostic under symmetric stochastic oscillations. In particular, signSGD can have $\E[w]\approx\theta$ while individual trajectories retain a large steady-state variance. In the current experiment, the final-500-sample signSGD mean signed error is almost zero, but its mean absolute error is about 0.205 and its RMSE about 0.256. The corresponding LIF-SGD ($\rho=0.98$) values are approximately 0.074 and 0.095.

\paragraph{Stale membrane evidence.}
Immediately before the switch, the IF membrane contains residual evidence primarily associated with motion toward the old optimum. At the switch, the mean IF membrane is about $-0.094$, and roughly 69\% of realizations have a negative membrane state, which opposes the new positive update direction after the optimum jumps to 2.

\paragraph{Oracle reset diagnostic.}
To isolate stale evidence, an oracle baseline resets the IF membrane to zero exactly at the objective change. The resulting performance change is very small: post-switch MSE over the first 1000 samples improves only from roughly 0.325 to 0.321, and the recovery time changes only by a few gradient evaluations. Therefore, for this large optimum shift, stale membrane evidence is not the dominant limitation. The new gradient is strong enough to reverse the accumulator rapidly.

\paragraph{Role of leak in this experiment.}
Increasing leakage mainly trades communication for responsiveness. Post-switch event counts decrease as $\rho$ decreases, but recovery becomes slower. The current evidence therefore does \emph{not} support the strong claim that leakage inherently improves adaptation to a changing objective.

\subsection{Experiment 3: moderate optimum shift after stationary convergence}

\paragraph{Hypothesis and design.}
A smaller objective shift was introduced to test the possibility that stale evidence matters when the new gradient is comparable to the residual membrane state. A first trial switching at $k=500$ was rejected as confounded: strongly leaky methods had not yet converged to the old optimum and therefore happened to start closer to the new optimum. The corrected experiment first lets all IF/LIF methods settle near the old optimum and then applies
\begin{equation}
  \theta_k=\begin{cases}
  0,&k<2000,\\
  0.5,&k\geq2000.
  \end{cases}
\end{equation}
The experiment uses 400 Monte Carlo runs. This long pre-switch phase makes the old-objective errors of SGD and IF/LIF comparable: the mean absolute errors immediately before the switch are approximately 0.078 for SGD, 0.077 for IF, 0.075 for $\rho=0.995$, 0.072 for $\rho=0.98$, and 0.066 for $\rho=0.95$.

\paragraph{Result.}
The moderate shift does not reveal a tracking advantage from leakage. Over the first 500 post-switch samples, the measured MAEs are approximately
\begin{center}
\begin{adjustbox}{max width=\textwidth}
\begin{tabular}{lrr}
\toprule
Method & post-switch MAE, $0{:}500$ & post-switch events \\
\midrule
SGD & 0.203 & 1500 \\
IF-SGD & 0.204 & 52.0 \\
IF-SGD, oracle reset & 0.204 & 51.9 \\
LIF-SGD, $\rho=0.995$ & 0.207 & 50.0 \\
LIF-SGD, $\rho=0.98$ & 0.217 & 43.4 \\
LIF-SGD, $\rho=0.95$ & 0.243 & 31.6 \\
\bottomrule
\end{tabular}
\end{adjustbox}
\end{center}
The oracle reset again changes the IF trajectory only marginally. Recovery to sustained mean absolute error below 0.1 takes about 397 gradient evaluations for IF and about 399 for oracle-reset IF, whereas stronger leakage delays recovery to roughly 421 evaluations at $\rho=0.98$ and 515 at $\rho=0.95$.

\paragraph{Interpretation.}
This experiment further weakens the initial hypothesis that leakage provides a direct adaptation advantage after an objective shift. Once pre-switch convergence is controlled, leak mainly reduces the rate at which new evidence is accumulated. The apparent benefit seen in the rejected early-switch trial was an initialization artifact. This is an important methodological lesson for later asynchronous experiments: methods must be compared at equivalent model states when the objective or global model changes.

\subsection{Experiment 4: controlled stale membrane}

\paragraph{Purpose.}
The fourth experiment removes the optimization trajectory entirely and isolates the stale-evidence mechanism. The membrane is initialized at a prescribed negative state
\begin{equation}
  z_0\in\{-0.45,-0.30,-0.10\},
\end{equation}
with threshold $\Delta=0.5$. At $k=0$, a constant new gradient $g=-G$, $G=0.4$, is applied, so the desired event direction is positive. With $\gamma=0.05$, the deterministic membrane input is
\begin{equation}
  u=\gamma G=0.02.
\end{equation}
The discrete dynamics are
\begin{equation}
  z_{k+1}=\rho z_k+u.
\end{equation}
Two first-passage quantities are distinguished:
\begin{enumerate}
  \item the time until stale sign is erased, i.e. $z_k\geq0$;
  \item the time until a useful positive event is emitted, i.e. $z_k\geq\Delta$.
\end{enumerate}

\paragraph{Exact deterministic expressions.}
For IF ($\rho=1$),
\begin{equation}
  k_0^{\mathrm{IF}}=\left\lceil\frac{-z_0}{u}\right\rceil,
  \qquad
  k_\Delta^{\mathrm{IF}}=\left\lceil\frac{\Delta-z_0}{u}\right\rceil.
\end{equation}
For LIF ($0<\rho<1$), define the forced equilibrium
\begin{equation}
  z_\infty=\frac{u}{1-\rho}.
\end{equation}
Then
\begin{equation}
  z_k=z_\infty+\rho^k(z_0-z_\infty).
\end{equation}
The stale sign is erased after
\begin{equation}
  k_0^{\mathrm{LIF}}
  =\left\lceil
  \frac{\log\!\left(z_\infty/(z_\infty-z_0)\right)}{\log\rho}
  \right\rceil.
\end{equation}
A positive communication event exists deterministically only if
\begin{equation}
  z_\infty>\Delta,
\end{equation}
which is equivalent to the discrete deadzone condition
\begin{equation}
  u>(1-\rho)\Delta.
\end{equation}
When this holds,
\begin{equation}
  k_\Delta^{\mathrm{LIF}}
  =\left\lceil
  \frac{\log\!\left((z_\infty-\Delta)/(z_\infty-z_0)\right)}{\log\rho}
  \right\rceil.
\end{equation}

\paragraph{Key result.}
Leak does indeed erase the stale sign faster. For $z_0=-0.45$, the deterministic zero-crossing times are 23, 22, 19, and 15 steps for $\rho=1$, 0.995, 0.98, and 0.95, respectively. However, the first useful threshold crossing is not accelerated: the corresponding $+\Delta$ passage times are 48, 48, 53, and no deterministic firing at all for $\rho=0.95$. The same qualitative behavior occurs for the other stale-state values.

This separates two effects that were previously conflated:
\begin{equation}
  \boxed{\text{faster forgetting of stale sign}\;\not\Rightarrow\;\text{faster useful communication}.}
\end{equation}
The reason is structural. While $z<0$, the leak term assists movement toward zero; after zero crossing, the same leak opposes accumulation toward $+\Delta$. A sufficiently large leak creates a deadzone in which weak but persistent new gradients cannot produce any deterministic event.

\paragraph{Noisy confirmation.}
A Monte Carlo test with gradient-noise standard deviation 0.4 confirms the deterministic picture. For $z_0=-0.45$, the mean zero-crossing time decreases from about 23.4 steps at $\rho=1$ to 19.1 at $\rho=0.98$, whereas the mean first-event time increases from about 48.4 to 53.2 steps. At $\rho=0.95$, stochastic noise eventually drives most runs across the threshold, but the mean delay rises to roughly 196 steps. Wrong first events are rare in this selected constant-gradient regime, so the dominant effect is the first-passage delay rather than event-sign reliability.

\paragraph{Implication for the role of leak.}
The controlled experiment falsifies the simple narrative that leak should be added because it makes an optimizer respond faster to changed gradients. Its stronger current justification is instead as a \emph{memory and communication regularizer}: it suppresses long-lived/noisy evidence and introduces an explicit freshness horizon, but the price is attenuated accumulation of weak persistent gradients and a possible communication deadzone. Any future asynchronous-FL argument for leakage must therefore demonstrate a net systems benefit under realistic model drift, not merely cite faster decay of stale membrane state.

\subsection{What these experiments establish---and what they do not}

The current evidence supports the following working observations:
\begin{enumerate}
  \item Event timing carries useful gradient-magnitude information.
  \item Temporal accumulation can dramatically improve signed-update reliability under noisy gradients.
  \item Sparse IF/LIF events can reproduce an SGD-like stochastic accuracy floor in the scalar diagnostic while using far fewer communication events.
  \item Leakage is a genuine memory/communication parameter, but stronger leak generally slows the accumulation of useful persistent evidence.
  \item Faster decay of stale membrane state does not by itself imply faster useful communication after a gradient reversal.
  \item Apparent nonstationary-tracking advantages can be artifacts if competing methods are not at comparable states when the objective changes.
\end{enumerate}

These experiments do not establish superiority over modern communication-efficient distributed optimization. They do not include vector event addresses, real bit accounting, client heterogeneity, partial participation, asynchronous server updates, or neural-network objectives. Their purpose is to validate and falsify mechanisms before those complications are introduced.

\section{Research Roadmap}

\subsection{Completed mechanism diagnostics}

The first eight experiments establish the following progression:
\begin{enumerate}
  \item \textbf{Stationary stochastic quadratic.} Temporal accumulation improves signed-event reliability and strongly reduces communication relative to periodic sign updates in the toy setting.
  \item \textbf{Large moving optimum.} Residual membrane evidence is present but is rapidly overwhelmed by a strong new gradient.
  \item \textbf{Moderate moving optimum after stationary convergence.} Leakage does not improve tracking and mainly reduces communication at the cost of slower adaptation.
  \item \textbf{Controlled stale membrane.} Leakage erases an opposing membrane sign faster but does not reduce the total first-passage time to a useful threshold event, and sufficiently strong leakage creates a deadzone.
  \item \textbf{Controlled asynchronous model drift.} The same distinction holds when the objective is fixed and the gradient changes only because the server model moves.
  \item \textbf{Endogenous two-client asynchronous learning.} A selected noisy operating point shows a net benefit for mild finite memory, initially suggesting soft stale-evidence invalidation.
  \item \textbf{Asynchrony--memory regime map.} The finite-memory optimum remains approximately $\rho=0.99$ for every tested compute-rate ratio, including $R=1$, and disappears entirely in a zero-noise homogeneous control. This falsifies the strong claim that the observed gain is primarily caused by asynchronous stale-evidence handling in the homogeneous problem.
  \item \textbf{Heterogeneous delayed-gradient federation with rate normalization.} True computation-delay staleness is introduced explicitly and can produce locally stale, globally harmful events under strong delay. However, a fresh-gradient oracle removes these stale events without materially improving IF, while LIF also improves in regimes where stale events are absent. Thus the dominant LIF gain in the scalar heterogeneous test is broader event regularization/deadzone behavior rather than a robust FL-specific freshness mechanism.
\end{enumerate}

The current interpretation of leak should therefore be explicit and conservative. The strongest supported roles are
\begin{equation}
  \boxed{\text{stochastic temporal-evidence filtering + event/communication regularization}.}
\end{equation}
Finite memory also creates a practical deadzone that can suppress persistent locally motivated but globally conflicting events under heterogeneous objectives. Whether that mechanism offers a genuine advantage over standard filtering, damping, deadband, proximal regularization, or error-feedback variants is now the central question.

\subsection{Working algorithmic reference: FedLIF-v0}

Before introducing larger benchmarks, the project should keep one simple reference formulation fixed. Client $i$ maintains a communication state $z_i$ and updates, coordinatewise,
\begin{equation}
  z_{i,j}^{k+1}=\rho z_{i,j}^{k}-\gamma_i g_{i,j}^{k},
\end{equation}
where $g_i^k$ is a local stochastic gradient. For asynchronous clients with intended aggregation weight $p_i$ and compute period $T_i$, the current rate-normalized scalar experiments use
\begin{equation}
  \gamma_i=\gamma p_iT_i,
\end{equation}
so unequal compute rates do not automatically redefine the global objective.

When
\begin{equation}
  |z_{i,j}|\geq\Delta_j,
\end{equation}
client $i$ emits a signed event
\begin{equation}
  s_{i,j}=\operatorname{sign}(z_{i,j})\in\{-1,+1\},
\end{equation}
and the global parameter receives the fixed jump
\begin{equation}
  w_j^+=w_j+q_js_{i,j}.
\end{equation}
A subtractive reset preserves threshold overshoot,
\begin{equation}
  z_{i,j}^+=z_{i,j}-\Delta_js_{i,j}.
\end{equation}
The three principal design variables remain deliberately separated:
\begin{equation}
  \boxed{\rho:\text{ memory},\qquad \Delta:\text{ evidence threshold},\qquad q:\text{ update quantum}.}
\end{equation}
This reference formulation is denoted \emph{FedLIF-v0} for planning purposes only; the name is not yet intended as a final algorithmic novelty claim.

\subsection{Planned experimental sequence}

The next experiments are ordered so that each one answers a specific unresolved question before additional complexity is introduced. A later experiment should only be pursued if the preceding stage leaves a technically meaningful mechanism worth carrying forward.

\subsubsection{Experiment 9: mechanism and baseline audit}

\paragraph{Question.}
Does the LIF mechanism provide a communication--accuracy trade-off that cannot be reproduced by simpler conventional filtering, residual-memory, or event-rejection mechanisms?

\paragraph{Why this experiment is necessary.}
Experiments~1--8 show that finite memory can improve event reliability, reduce communication, and suppress persistent heterogeneous event activity. However, these observations are not yet uniquely neuromorphic. Similar effects may be obtained from exponential moving-average filtering, memory-decayed error feedback, server-side deadbands, or simple reset rules. Scaling to neural networks before resolving this question risks building a large benchmark around a mechanism that is only a reformulation of known ideas.

\paragraph{Setup.}
Use the current heterogeneous asynchronous scalar quadratic and, where needed, the delayed-gradient variant. Compare FedLIF-v0 against at least:
\begin{itemize}
  \item periodic signSGD;
  \item IF with $\rho=1$;
  \item exponential moving-average gradient filtering followed by threshold/sign communication;
  \item decayed error-feedback or residual-memory compression;
  \item hard and partial membrane reset after remote updates;
  \item server-side deadband/event rejection;
  \item a proximal local update or another simple regularizer that suppresses persistent local disagreement.
\end{itemize}
Each method must receive its own parameter sweep rather than being compared at one hand-selected operating point.

\paragraph{Hypothesis.}
If temporal evidence accumulation plus thresholded event generation is genuinely useful, FedLIF should achieve a non-dominated region in the communication--accuracy plane rather than merely matching one of the simpler alternatives.

\paragraph{Primary metrics.}
The primary comparison is a Pareto front of
\begin{equation}
  \text{optimization error}\quad\text{versus}\quad\text{transmitted bits}.
\end{equation}
Secondary metrics are event count, locally wrong/stale events, globally harmful events, per-client participation, and whole-run convergence cost.

\paragraph{Decision criterion.}
Proceed with the LIF mechanism only if FedLIF has a robust non-dominated region relative to the closest conventional comparators. If a simpler EMA, deadband, or decayed-error-feedback baseline reproduces the same Pareto front, the algorithmic novelty must be reconsidered before scaling.

\subsubsection{Experiment 10: vector heterogeneous asynchronous quadratics}

\paragraph{Question.}
Does the mechanism survive when communication is genuinely high-dimensional and heterogeneous gradients disagree across coordinates rather than only along one scalar direction?

\paragraph{Why this experiment is necessary.}
The scalar setting hides several important effects: coordinate addressing is free, simultaneous threshold crossings cannot occur, curvature is trivial, and a signed event has no spatial structure. A communication-efficient federated optimizer must work when $w\in\mathbb{R}^d$ and when clients disagree in multiple directions.

\paragraph{Setup.}
Use $N$ heterogeneous clients with strongly convex local quadratics
\begin{equation}
  F_i(w)=\frac{1}{2}(w-\theta_i)^\top H_i(w-\theta_i),
\end{equation}
with different $\theta_i$, positive-definite $H_i$, stochastic gradients, heterogeneous compute rates, and asynchronous completion delays. Start with $d\in[20,50]$ and $N\approx10$. Because the aggregate optimum is known analytically,
\begin{equation}
  w^\star=\left(\sum_i p_iH_i\right)^{-1}\sum_i p_iH_i\theta_i,
\end{equation}
optimization error can still be measured exactly.

\paragraph{Additional communication question.}
An event now requires an address as well as a sign. A coordinate event therefore costs approximately
\begin{equation}
  B_{\mathrm{event}}\approx\lceil\log_2 d\rceil+1
\end{equation}
logical bits before packet/protocol overhead. The experiment must therefore count actual encoded bits rather than treating a signed event as ``one bit.''

\paragraph{Hypothesis.}
If the mechanism is scalable, temporal evidence should still concentrate communication on persistently informative coordinates and generate a favorable error--bit trade-off relative to periodic sign and sparse/error-feedback baselines.

\paragraph{Primary metrics.}
Report error versus bits, error versus event count, per-client and per-coordinate firing rates, client-weight fidelity, and the distribution of coordinate participation.

\paragraph{Decision criterion.}
The method should retain a communication advantage without silencing slow clients, distorting the intended weights $p_i$, or requiring an event address cost that eliminates the benefit.

\subsubsection{Experiment 11: convergence refinement and adaptive resolution}

\paragraph{Question.}
Can the event-driven optimizer retain sparse communication while avoiding the permanent practical-convergence neighborhood induced by fixed update quantum $q$ and finite-memory deadzones?

\paragraph{Why this experiment is necessary.}
The scalar experiments repeatedly show that fixed $q$ and fixed $\Delta$ naturally lead to practical rather than exact convergence. A serious optimizer needs a mechanism for coarse early motion and fine late optimization.

\paragraph{Setup.}
Starting from the vector quadratic, compare:
\begin{itemize}
  \item fixed $q$ and fixed $\Delta$;
  \item decaying $q_k$ with fixed $\Delta$;
  \item coupled schedules $q_k=c\Delta_k$;
  \item adaptive thresholds $\Delta_k$;
  \item optional homeostatic threshold adaptation that targets a desired event rate.
\end{itemize}
A generic decreasing update quantum may take the form
\begin{equation}
  q_k=\frac{q_0}{(1+k/k_0)^\alpha}.
\end{equation}

\paragraph{Hypothesis.}
A multiscale event representation should allow large early updates with high event activity, followed by progressively finer and sparser updates near the optimum.

\paragraph{Primary metrics.}
Measure asymptotic optimization error, convergence rate, total bits, event-rate trajectory, and sensitivity to the schedule parameters.

\paragraph{Decision criterion.}
The refined method should reduce or remove the fixed-$q$ accuracy floor without losing the communication advantage established in Experiment~10.

\subsubsection{Experiment 12: federated logistic regression}

\paragraph{Question.}
Does the event-driven mechanism remain useful on a real high-dimensional stochastic learning problem before introducing nonconvex neural networks?

\paragraph{Why this experiment is necessary.}
Logistic regression provides a real data-distributed learning problem with high-dimensional stochastic gradients while retaining a comparatively interpretable convex objective. It is therefore the appropriate bridge between synthetic vector quadratics and neural-network training.

\paragraph{Setup.}
Use approximately $N=20$ clients with non-IID data partitions, e.g., Dirichlet label heterogeneity, heterogeneous compute rates, and asynchronous updates. Candidate datasets include a binary or multiclass MNIST-derived task or another compact classification benchmark. Compare FedLIF against the mechanism winners from Experiment~9 plus standard local-SGD/FedAvg, periodic sign communication, sparse ternary/error-feedback compression, and an event-triggered/lazy distributed baseline.

\paragraph{Hypothesis.}
FedLIF should convert persistent stochastic learning signal into an adaptive event rate, communicating frequently during informative phases and becoming quieter as useful evidence decreases.

\paragraph{Primary metrics.}
The central result must be
\begin{equation}
  \boxed{\text{test loss/accuracy versus cumulative transmitted bits}.}
\end{equation}
Secondary axes are wall-clock simulation time, gradient evaluations, communication events, per-client participation, and event-rate evolution over training.

\paragraph{Decision criterion.}
Proceed to a neural network only if the event-driven method remains competitive in accuracy and has a meaningful communication advantage after realistic address and packet costs are included.

\subsubsection{Experiment 13: small neural-network benchmark}

\paragraph{Question.}
Can neuromorphic temporal evidence coding train an ordinary nonconvex neural network under heterogeneous federated data while remaining communication efficient?

\paragraph{Why this experiment is necessary.}
The ultimate target is not an SNN. Backpropagation should remain conventional; the neuromorphic contribution is the communication/optimization layer after gradient computation. This experiment tests whether the same mechanism survives realistic high-dimensional, noisy, nonconvex gradients.

\paragraph{Setup.}
Begin with a small MLP on Fashion-MNIST or MNIST. Only after that is stable should a modest CNN/CIFAR-10 setting be considered. Coordinatewise LIF should be compared with blockwise/grouped event encoders because raw coordinate addresses may dominate communication cost in large networks.

\paragraph{Hypothesis.}
The characteristic neuromorphic behavior should remain visible:
\begin{equation}
  \text{high event activity early in training}\quad\rightarrow\quad\text{sparser firing later},
\end{equation}
with gradient magnitude represented partly through event timing/rate rather than transmitted amplitude.

\paragraph{Primary metrics.}
Report accuracy/loss versus transmitted bits, event-rate trajectories, firing sparsity by layer/block, client participation, and comparison against the strongest compression/error-feedback baselines retained from earlier experiments.

\paragraph{Decision criterion.}
A neural-network result is useful only if the communication advantage survives realistic encoding and does not rely on silently suppressing difficult or slow clients.

\subsection{Cross-cutting communication architecture}

The eventual architecture should distinguish between \emph{event generation} and \emph{network packetization}. A neuromorphic encoder may generate asynchronous coordinate events internally, but practical communication should batch events over a short interval,
\begin{equation}
  \mathcal{E}_i=\{(j_1,s_1),\ldots,(j_m,s_m)\},
\end{equation}
and transmit them using an efficient sparse representation. Possible encodings include sorted indices with delta coding, bitmaps, or block addresses. The bit accounting must include index, sign, headers, and any synchronization/checkpoint overhead.

The same event representation may eventually support the downlink. If the server model evolves only through accepted jumps
\begin{equation}
  w_j\leftarrow w_j+q_js,
\end{equation}
then clients can reconstruct the global model from the broadcast event stream after an initial synchronization. This creates the possibility of a fully event-coded uplink and downlink, although periodic checkpoints will likely be required for robustness.

\subsection{Theory milestones}

The theoretical progression should follow the mechanism that is actually supported by the experiments:
\begin{enumerate}
  \item deterministic scalar hybrid model and practical-convergence bound;
  \item stochastic scalar first-passage analysis and expected descent;
  \item explicit noise--memory design rule for the LIF evidence state;
  \item characterization of the finite-memory deadzone and event-suppression behavior under heterogeneous local gradients;
  \item comparison/equivalence results between LIF memory, exponential residual decay, error-feedback memory decay, EMA filtering, and simpler deadband mechanisms;
  \item vector strongly convex objectives with coordinatewise or blockwise event encoders;
  \item diminishing/adaptive update quanta and thresholds for asymptotic convergence;
  \item stochastic smooth objectives with bounded variance and asynchronous updates;
  \item multi-client objectives $F=\sum_i p_iF_i$ with bounded heterogeneity, computation-rate mismatch, and communication delay;
  \item nonconvex neural-network settings.
\end{enumerate}

The theoretical quantities should mirror the control perspective where useful: Lyapunov/objective decrease, practical convergence neighborhoods, event-rate bounds, robustness to noisy inputs, communication-deadzone conditions, and the relationship between evidence memory and local gradient statistics. The previously proposed direct relation
\begin{equation}
  \tau_{\mathrm{mem}}\sim\tau_{\mathrm{remote\ model\ change}}
\end{equation}
should not be assumed; Experiments~7--8 do not support it as a general explanation of the observed finite-memory benefit.

\subsection{Paper-level decision criteria}

The project should proceed toward a paper only if at least one of the following stronger claims survives systematic comparison with existing methods:
\begin{itemize}
  \item LIF/IF temporal evidence gives a better communication--reliability trade-off than periodic sign and standard error-feedback compression under stochastic gradients;
  \item the hybrid-system formulation yields a useful noise--memory, threshold, or event-rate design law that is not obvious from standard compression analysis;
  \item finite-memory event suppression under heterogeneous clients outperforms simpler deadband, proximal, EMA, reset, or decayed-error-feedback mechanisms at matched communication and accuracy;
  \item asynchronous event timing provides a systems-level benefit under heterogeneous client computation and communication rates without silencing slow clients or distorting intended client weights;
  \item adaptive/homeostatic firing thresholds create a robust communication-budget mechanism that conventional thresholding does not reproduce as cleanly;
  \item bidirectional event-coded model evolution materially reduces communication while preserving useful learning behavior.
\end{itemize}

A direct claim that LIF is valuable because it fixes stale asynchronous FL gradients is no longer supported by the current scalar evidence and should not be used as the core paper story unless a later, substantially different setting overturns Experiments~7--8.

If none of the remaining claims survives modern baselines and realistic bit accounting, the project should be reframed as a hybrid-systems interpretation of existing error-feedback/event-triggered optimization rather than forced into a novelty claim.

\subsection{Working name}

The repository and umbrella project are named \texttt{NeuromorphicFL}. The reference algorithm is provisionally denoted \emph{FedLIF-v0} in the roadmap. A final algorithm name should only be fixed after Experiment~9 establishes whether the mechanism is actually distinct from the closest conventional alternatives.


\bibliographystyle{unsrtnat}
\bibliography{references}

\end{document}
